\documentclass[10pt,a4paper]{amsart}
\usepackage{amsmath,amssymb,amsthm,graphicx}
% Retained preamble for the Stacks Project chapter snapshot in this directory.
% The chapter's own preamble.tex is a copy of the shared Stacks Project preamble; the block
% below is sliced MECHANICALLY, line for line, from it. Nothing was added, removed or
% redefined; the wrapper class, the xr-hyper/hyperref wiring and the externaldocument list
% are omitted because the assembler supplies the document class.
%
% lines taken from preamble.tex: 3, 5, 9, 11, 12, 156, 157, 158, ... (51 lines total)
%
% package coverage restored from the chapter's own preamble (2026-10-13): xy, xy-2cell

\usepackage[all]{xy}
\xyoption{2cell}
\UseAllTwocells

\newenvironment{reference}{\comment}{\endcomment}
\newenvironment{slogan}{\comment}{\endcomment}
\newenvironment{history}{\comment}{\endcomment}
\theoremstyle{plain}
\newtheorem{theorem}[subsection]{Theorem}
\newtheorem{proposition}[subsection]{Proposition}
\newtheorem{lemma}[subsection]{Lemma}
\theoremstyle{definition}
\newtheorem{definition}[subsection]{Definition}
\newtheorem{example}[subsection]{Example}
\newtheorem{exercise}[subsection]{Exercise}
\newtheorem{situation}[subsection]{Situation}
\theoremstyle{remark}
\newtheorem{remark}[subsection]{Remark}
\newtheorem{remarks}[subsection]{Remarks}
\numberwithin{equation}{subsection}
\def\lim{\mathop{\mathrm{lim}}\nolimits}
\def\colim{\mathop{\mathrm{colim}}\nolimits}
\def\Spec{\mathop{\mathrm{Spec}}}
\def\Hom{\mathop{\mathrm{Hom}}\nolimits}
\def\Ext{\mathop{\mathrm{Ext}}\nolimits}
\def\SheafHom{\mathop{\mathcal{H}\!\mathit{om}}\nolimits}
\def\SheafExt{\mathop{\mathcal{E}\!\mathit{xt}}\nolimits}
\def\Sch{\mathit{Sch}}
\def\Mor{\mathop{\mathrm{Mor}}\nolimits}
\def\Ob{\mathop{\mathrm{Ob}}\nolimits}
\def\Sh{\mathop{\mathit{Sh}}\nolimits}
\def\NL{\mathop{N\!L}\nolimits}
\def\CH{\mathop{\mathrm{CH}}\nolimits}
\def\proetale{{pro\text{-}\acute{e}tale}}
\def\etale{{\acute{e}tale}}
\def\QCoh{\mathit{QCoh}}
\def\Ker{\mathop{\mathrm{Ker}}}
\def\Im{\mathop{\mathrm{Im}}}
\def\Coker{\mathop{\mathrm{Coker}}}
\def\Coim{\mathop{\mathrm{Coim}}}
\DeclareMathSymbol{\boxtimes}{\mathbin}{AMSa}{"02}
\def\QCohstack{\mathcal{QC}\!\mathit{oh}}
\def\Cohstack{\mathcal{C}\!\mathit{oh}}
\def\Spacesstack{\mathcal{S}\!\mathit{paces}}
\def\Quotfunctor{\mathrm{Quot}}
\def\Hilbfunctor{\mathrm{Hilb}}
\def\Curvesstack{\mathcal{C}\!\mathit{urves}}
\def\Polarizedstack{\mathcal{P}\!\mathit{olarized}}
\def\Complexesstack{\mathcal{C}\!\mathit{omplexes}}
\def\Pic{\mathop{\mathrm{Pic}}\nolimits}
\def\Picardstack{\mathcal{P}\!\mathit{ic}}
\def\Picardfunctor{\mathrm{Pic}}
\def\Deformationcategory{\mathcal{D}\!\mathit{ef}}
\usepackage{cleveref}
\crefname{originallemma}{original lemma}{original lemmas}
\Crefname{originallemma}{Original lemma}{Original lemmas}
\makeatletter
\newlabel{simplicial-lemma-cosk-shriek}{{018V}{source}{}{sourcepointer-1}{}}
\newlabel{simplicial-lemma-cosk-shriek@cref}{{[originallemma][0][]018V}{[1][0][]source}{}{}{}}
\makeatother
\crefname{originallemma}{original lemma}{original lemmas}
\Crefname{originallemma}{Original lemma}{Original lemmas}
\makeatletter
\newlabel{simplicial-lemma-exists-hom-from-simplicial-set-finite}{{017L}{source}{}{sourcepointer-2}{}}
\newlabel{simplicial-lemma-exists-hom-from-simplicial-set-finite@cref}{{[originallemma][0][]017L}{[1][0][]source}{}{}{}}
\makeatother
\crefname{originallemma}{original lemma}{original lemmas}
\Crefname{originallemma}{Original lemma}{Original lemmas}
\makeatletter
\newlabel{simplicial-lemma-hom-from-coprod}{{017N}{source}{}{sourcepointer-3}{}}
\newlabel{simplicial-lemma-hom-from-coprod@cref}{{[originallemma][0][]017N}{[1][0][]source}{}{}{}}
\makeatother
\title{Identifying the terms of the Hom(Delta[1], L) coskeleton diagram}
\author{The Stacks Project Authors}
\date{stacks-project:hypercovering.tex:local-snapshot-20261012; GNU FDL 1.2 or later}
\begin{document}
\maketitle
\noindent\textit{Editorial note.} Statement span is lines 1995-2054 of sources/stacks-hypercovering/pinned-original-source.tex; proof span is lines 2056-2072, the complete author proof (15 body lines).\par\medskip
\noindent\textit{Editorial note.} Navigation stubs: 3 cited-but-not-included labels are supplied as non-load-bearing pointers with their pinned tags where the snapshot tag index has them.\par\medskip
\noindent\textit{Editorial note.} extra\_packages is empty on purpose: the retained preamble already loads xy with options and a second declaration clashes in the assembler wrapper.\par\medskip
\noindent\textit{Editorial note.} Text is retained verbatim, including the author's own typos and placeholders.\par\medskip
\noindent\textit{Editorial note.} Queue index 20 of queue\_hypercovering.json (pb=745, stmt 1995-2054, proof 2056-2072).\par\medskip
\noindent\textit{Editorial note.} stacks-hypercovering wave-1 (hc01): fresh adjudication accept (2026-10-14); pool per standing fallback order (hypercovering -> pic -> spaces-more-cohomology) after md01 closed the moduli pool; hypercovering.tex 53 labels/34 with proof judged 26 accept / 8 reject / 14 pb=0 env-dup rows excluded; spans cite-free and reference/slogan/history-free so no patchenv needed. Named omission: The intersection identity (partial Delta[n+1] coprod partial Delta[n+1] = the intersection inside Delta[1] x Delta[n+1]) is 'left to the reader' - a named mechanical check; the four identifications and the pushout argument are in-span.\par\medskip
\noindent\textit{Allowed standard prerequisites (editorial).}\begin{itemize}
\item Explicit prerequisite boundary. The cited inputs are prerequisites; the argument in the retained span is not replaced by any cited result.
\end{itemize}\medskip
\ifdefined\hypertarget\hypertarget{sourcepointer-1}{}\fi
\noindent\textit{Source reference.} Original lemma 018V: sources/stacks-hypercovering/pinned-original-source.tex, source lines 2072--2072. Cited by the retained statement or proof but not part of the collected spans; supplied only as a navigation pointer. The argument does not depend on its content.\par\medskip
\ifdefined\hypertarget\hypertarget{sourcepointer-2}{}\fi
\noindent\textit{Source reference.} Original lemma 017L: sources/stacks-hypercovering/pinned-original-source.tex, source lines 2072--2072. Cited by the retained statement or proof but not part of the collected spans; supplied only as a navigation pointer. The argument does not depend on its content.\par\medskip
\ifdefined\hypertarget\hypertarget{sourcepointer-3}{}\fi
\noindent\textit{Source reference.} Original lemma 017N: sources/stacks-hypercovering/pinned-original-source.tex, source lines 2072--2072. Cited by the retained statement or proof but not part of the collected spans; supplied only as a navigation pointer. The argument does not depend on its content.\par\medskip
\begin{lemma}
\label{lemma-hom-hypercovering}
Let $\mathcal{C}$ be a site with fibre products.
Let $X$ be an object of $\mathcal{C}$.
Let $L$ be a simplicial object of $\text{SR}(\mathcal{C}, X)$.
Let $n \geq 0$. Consider the commutative diagram
\begin{equation}
\label{equation-diagram}
\xymatrix{
\Hom(\Delta[1], L)_{n + 1} \ar[r] \ar[d] &
(\text{cosk}_n \text{sk}_n \Hom(\Delta[1], L))_{n + 1} \ar[d] \\
(L \times L)_{n + 1} \ar[r] &
(\text{cosk}_n \text{sk}_n (L \times L))_{n + 1}
}
\end{equation}
coming from the morphism defined above.
We can identify the terms in this diagram as follows,
where
$\partial \Delta[n + 1] = i_{n!}\text{sk}_n \Delta[n + 1]$
is the $n$-skeleton of the $(n + 1)$-simplex:
\begin{eqnarray*}
\Hom(\Delta[1], L)_{n + 1}
& = &
\Hom(\Delta[1] \times \Delta[n + 1], L)_0 \\
(\text{cosk}_n \text{sk}_n \Hom(\Delta[1], L))_{n + 1}
& = &
\Hom(\Delta[1] \times \partial \Delta[n + 1], L)_0 \\
(L \times L)_{n + 1}
& = &
\Hom(
(\Delta[n + 1] \amalg \Delta[n + 1], L)_0 \\
(\text{cosk}_n \text{sk}_n (L \times L))_{n + 1}
& = &
\Hom(
\partial \Delta[n + 1]
\amalg
\partial \Delta[n + 1], L)_0
\end{eqnarray*}
and the morphism between these objects of $\text{SR}(\mathcal{C}, X)$
come from the commutative diagram of simplicial sets
\begin{equation}
\label{equation-dual-diagram}
\xymatrix{
\Delta[1] \times \Delta[n + 1] &
\Delta[1] \times \partial\Delta[n + 1] \ar[l] \\
\Delta[n + 1] \amalg \Delta[n + 1] \ar[u] &
\partial\Delta[n + 1] \amalg \partial\Delta[n + 1]
\ar[l] \ar[u]
}
\end{equation}
Moreover the fibre product of the bottom arrow and the
right arrow in (\ref{equation-diagram}) is equal to
$$
\Hom(U, L)_0
$$
where $U \subset \Delta[1] \times \Delta[n + 1]$
is the smallest simplicial subset such that both
$\Delta[n + 1] \amalg \Delta[n + 1]$ and
$\Delta[1] \times \partial\Delta[n + 1]$ map into it.
\end{lemma}

\begin{proof}
The first and third equalities are
Simplicial, Lemma \ref{simplicial-lemma-exists-hom-from-simplicial-set-finite}.
The second and fourth follow from the cited lemma combined with
Simplicial, Lemma \ref{simplicial-lemma-cosk-shriek}.
The last assertion follows from the fact that
$U$ is the push-out of the bottom and right arrow of the
diagram (\ref{equation-dual-diagram}), via
Simplicial, Lemma \ref{simplicial-lemma-hom-from-coprod}.
To see that $U$ is equal to this push-out it suffices
to see that the intersection of
$\Delta[n + 1] \amalg \Delta[n + 1]$ and
$\Delta[1] \times \partial\Delta[n + 1]$
in $\Delta[1] \times \Delta[n + 1]$ is equal to
$\partial\Delta[n + 1] \amalg \partial\Delta[n + 1]$.
This we leave to the reader.
\end{proof}

\end{document}
